% !TEX program = xelatex
\documentclass[aspectratio=169,10pt]{ctexbeamer}

% ============ Theme & Packages ============
\usetheme{Madrid}
\usecolortheme{seahorse}
\setbeamertemplate{navigation symbols}{}
\setbeamertemplate{caption}[numbered]
\setbeamertemplate{frametitle continuation}{}

\usepackage{amsmath,amssymb,mathtools,bm,bbm}
\usepackage{booktabs}
\usepackage{tikz}
\usepackage{xcolor}
\usetikzlibrary{arrows.meta,positioning,shapes.geometric,calc,decorations.pathreplacing}

% ============ Math Notation ============
\newcommand{\R}{\mathbb{R}}
\newcommand{\E}{\mathbb{E}}
\newcommand{\diag}{\operatorname{diag}}
\newcommand{\tr}{\operatorname{tr}}
\newcommand{\vect}{\operatorname{vec}}
\newcommand{\eps}{\varepsilon}
\newcommand{\grad}{\nabla}
\newcommand{\norm}[1]{\left\lVert #1 \right\rVert}
\newcommand{\ip}[2]{\left\langle #1,#2 \right\rangle}
\newcommand{\Id}{\mathbb{I}}
\newcommand{\calN}{\mathcal{N}}
\newcommand{\calO}{\mathcal{O}}
\newcommand{\calL}{\mathcal{L}}
\newcommand{\calF}{\mathcal{F}}
\DeclareMathOperator*{\argmin}{arg\,min}
\DeclareMathOperator*{\argmax}{arg\,max}
\DeclareMathOperator{\sgn}{sgn}
\DeclareMathOperator{\Proj}{Proj}
\DeclareMathOperator{\Dist}{Dist}

\title[优化器与复杂度]{神经网络优化算法：\\[0.3em] 从算法复杂度到 AdamW、Muon 与 PolarGrad}
\subtitle{融合随机优化理论、自适应方法、与矩阵结构优化器}
\author{基于 Qi Deng \& 苏炜杰等讲义整理}
\date{\today}

\begin{document}

\begin{frame}
  \titlepage
\end{frame}

% ======================================================================
\section{引言与框架}
% ======================================================================

\begin{frame}{本文档的定位}
这是一个\textbf{学术级讲义}，融合：
\begin{itemize}
  \item \textbf{随机优化理论}（SAA、SA/SGD、Mirror Descent、方差缩减）——理解算法复杂度的数学语言
  \item \textbf{深度学习优化器}（Adam/AdamW、Shampoo、SOAP、Muon、PolarGrad）——实践中的算法设计
  \item \textbf{数学基础补充}（谱分解、SVD、极分解、矩阵范数、Newton-Schulz、浓度不等式）——让你能读懂论文
\end{itemize}

\vspace{0.3cm}
\textbf{阅读建议}：
\begin{itemize}
  \item 如果你只关心「AdamW vs Muon 用哪个」，先读 \S6--\S9
  \item 如果你想知道「为什么 Muon 有效」，请读 \S10（最陡下降统一视角）和 \S11（矩阵函数）
  \item 如果你想做理论科研，\S2--\S5（复杂度分析）是必修的
\end{itemize}
\end{frame}

\begin{frame}{为什么矩阵结构对优化器设计至关重要？}
神经网络参数天然是\textbf{矩阵}（权重层 $W\in\R^{m\times n}$）或\textbf{张量}（卷积核）。

\vspace{0.2cm}
\begin{columns}[T]
\column{0.45\textwidth}
\textbf{向量视角（Adam 的做法）}
\begin{itemize}
  \item 将 $W$ 展平为 $\vect(W)\in\R^{mn}$
  \item 逐坐标独立地估计二阶矩
  \item 忽略 $m$ 个行与 $n$ 个列之间的结构
  \item 预条件器是对角阵：$P=\diag(v_t)^{-1/2}$
\end{itemize}

\column{0.45\textwidth}
\textbf{矩阵视角（Muon/Shampoo 的做法）}
\begin{itemize}
  \item 保持 $W$ 的矩阵结构
  \item 分别处理行空间与列空间
  \item 利用梯度矩阵 $G_t$ 的 SVD/极分解
  \item 预条件器利用 Kronecker 结构或谱正交化
\end{itemize}
\end{columns}

\vspace{0.3cm}
\begin{center}
\begin{tikzpicture}[>=Stealth]
  \node[draw, rounded corners, fill=blue!8, minimum width=2.5cm, align=center] (v) {向量化\\$\vect(W)$};
  \node[draw, rounded corners, fill=red!8, minimum width=2.5cm, right=2cm of v, align=center] (m) {矩阵保持\\$W$};
  \node[below=0.8cm of v, font=\footnotesize] {Adam, AdaGrad};
  \node[below=0.8cm of m, font=\footnotesize] {Shampoo, Muon, SOAP};
  \draw[->] (v) -- (m) node[midway, above, font=\small] {\textbf{关键转变}};
\end{tikzpicture}
\end{center}
\end{frame}

% ======================================================================
\section{算法复杂度基础：随机优化}
% ======================================================================

\begin{frame}{随机优化问题的两种范式}
\textbf{问题设定}：$\min_{x\in\mathcal{X}} f(x) = \E_\xi[F(x,\xi)]$，其中 $\xi$ 是随机变量。

\vspace{0.3cm}
\begin{block}{范式一：样本平均近似 (SAA)}
  \begin{itemize}
    \item 一次性采样 $N$ 个 $\xi_i$，构造 $\hat f_N(x)=\frac1N\sum_{i=1}^N F(x,\xi_i)$
    \item 用确定性优化方法求解 $\min_x \hat f_N(x)$
    \item \textbf{核心问题}：$N$ 取多大才能保证 $\hat f_N$ 的解 $\varepsilon$-接近 $f$ 的真解？
  \end{itemize}
\end{block}

\vspace{0.2cm}
\begin{block}{范式二：随机近似 (SA / SGD)}
  \begin{itemize}
    \item 每步在线采样 $\xi_t$，获得随机梯度 $g_t = \grad_x F(x_t,\xi_t)$
    \item 迭代更新：$x_{t+1} = x_t - \eta_t g_t$
    \item \textbf{核心问题}：需要多少步 $T$ 才能收敛？步长 $\eta_t$ 如何选？
  \end{itemize}
\end{block}
\end{frame}

\begin{frame}{SAA 的样本复杂度}
\textbf{核心结论}：对有限域 $|\mathcal{X}|<+\infty$，若每个 $F(\cdot,\xi)$ 的波动是 sub-Gaussian（宽度 $\sigma$），则：

\begin{theorem}[SAA 样本复杂度]
  对任意 $0<\hat\varepsilon<\varepsilon$ 和置信度 $1-\delta$，若样本量
  \[
    N \ge \frac{2\sigma^2}{(\varepsilon-\hat\varepsilon)^2}\ln\frac{|\mathcal{X}|}{\delta},
  \]
  则 $\Pr(\hat{\mathcal{S}}_N^{\hat\varepsilon}\subseteq \mathcal{S}^\varepsilon)\ge 1-\delta$，
  即 SAA 的近似解集以高概率包含在真实问题的近似解集中。
\end{theorem}

\vspace{0.2cm}
\textbf{证明工具链}：
\[
  \text{Chebyshev} \;\to\; \text{MGF bound} \;\to\; \text{Hoeffding 不等式}
\]
对每个 $x\in\mathcal{X}\setminus\mathcal{S}^\varepsilon$，联合界给出上述结果。

\textbf{对连续域}：需引入覆盖数（covering number），复杂度变为
$\mathcal{O}\!\left(\frac{\sigma^2}{(\varepsilon-\hat\varepsilon)^2}\left[d\ln\frac{LD}{\varepsilon-\hat\varepsilon}+\ln\frac{1}{\delta}\right]\right)$。
\end{frame}

\begin{frame}{Sub-Gaussian 随机变量与浓度不等式（知识补充）}
\begin{definition}[Sub-Gaussian 随机变量]
  称零均值随机变量 $Y$ 是 $\sigma$-sub-Gaussian，若其矩母函数满足：
  \[
    \E[e^{tY}] \le \exp\!\left(\frac{\sigma^2 t^2}{2}\right),\quad \forall t\in\R.
  \]
\end{definition}

\vspace{0.2cm}
\textbf{常见 Sub-Gaussian 分布}：
\begin{itemize}
  \item $Y\sim\calN(0,\sigma^2)$：高斯（宽度 $\sigma$）
  \item $Y\in[a,b]$ 有界：Hoeffding 引理给出宽度 $(b-a)/2$
  \item Bernoulli($p$)：宽度受控
\end{itemize}

\vspace{0.2cm}
\begin{theorem}[Hoeffding 不等式]
  若 $Y_1,\ldots,Y_N$ 是独立 $\sigma$-sub-Gaussian 随机变量，则
  \[
    \Pr\!\left(\frac{1}{N}\sum_{i=1}^N Y_i \ge t\right) \le \exp\!\left(-\frac{N t^2}{2\sigma^2}\right).
  \]
\end{theorem}

\textbf{直观}：经验均值偏离真实期望超过 $t$ 的概率按 $\exp(-N t^2)$ 衰减——样本越多，浓度越紧。
\end{frame}

\begin{frame}{随机梯度下降的收敛保证}
对光滑凸函数 $f$，取 SGD 步长 $\eta_t = \frac{D}{\sigma\sqrt{T}}$（常数）或 $\eta_t = \frac{1}{L+t}$（递减），Polyak 平均 $\bar x_T = \frac1T\sum_t x_t$ 满足：
\[
  \E[f(\bar x_T)] - f(x^\star) \le \frac{D^2}{2\eta T} + \frac{\eta\sigma^2}{2}
  = \mathcal{O}\!\left(\frac{D\sigma}{\sqrt{T}}\right).
\]

\vspace{0.2cm}
与 GD 对比：
\begin{center}
\begin{tabular}{lcc}
\toprule
方法 & 每步开销 & 收敛速度 \\
\midrule
GD（全梯度） & $\mathcal{O}(n)$ & $\mathcal{O}(1/T)$（凸）、$\mathcal{O}(e^{-T/\kappa})$（强凸）\\
SGD（随机） & $\mathcal{O}(1)$ & $\mathcal{O}(1/\sqrt{T})$（凸）、$\mathcal{O}(1/T)$（强凸+递减步长）\\
\bottomrule
\end{tabular}
\end{center}

\textbf{关键的方差项}：$\frac{\eta\sigma^2}{2}$ 不随 $T$ 消失。这就是为什么\textbf{方差缩减}（variance reduction）和\textbf{自适应步长}（adaptive methods）是随机优化的两大主题。
\end{frame}

% ======================================================================
\section{深度学习训练的基础策略}
% ======================================================================

\begin{frame}{深度学习优化的特殊挑战}
\begin{columns}[T]
\column{0.5\textwidth}
\begin{enumerate}
  \item \textbf{非凸性}：损失曲面有大量局部极小和鞍点
  \item \textbf{高维性}：参数量 $10^6$--$10^{12}$
  \item \textbf{病态性}：不同参数方向的曲率差异极大（$\kappa\gg1$）
  \item \textbf{随机噪声}：mini-batch 引入方差
\end{enumerate}

\column{0.5\textwidth}
\begin{enumerate}
  \setcounter{enumi}{4}
  \item \textbf{梯度消失/爆炸}：深层网络梯度传播问题
  \item \textbf{过拟合}：模型记忆训练数据
  \item \textbf{超参数敏感}：学习率、batch size 需精细调节
  \item \textbf{计算约束}：GPU 内存、通信带宽
\end{enumerate}
\end{columns}

\vspace{0.3cm}
\textbf{这些挑战直接驱动了优化器设计的演进}：
\[
  \text{病态性 + 高维性} \;\to\; \text{预条件方法}\qquad
  \text{随机噪声} \;\to\; \text{EMA/动量}\qquad
  \text{计算约束} \;\to\; \text{内存高效的近似}
\]
\end{frame}

\begin{frame}{Weight Decay 与 L2 正则化的关键区别}
\textbf{L2 正则化}修改目标函数：
\[
  \mathcal{L}_{\text{L2}}(w) = \mathcal{L}(w) + \frac{\lambda}{2}\norm{w}^2,\quad
  g_t^{\text{L2}} = \grad\mathcal{L}(w_t) + \lambda w_t.
\]
在 SGD 中等价于 weight decay：$w_{t+1} = (1-\eta\lambda)w_t - \eta\grad\mathcal{L}(w_t)$。

\vspace{0.2cm}
\textbf{但在 Adam 中不等价！} 因为自适应缩放 $\frac{1}{\sqrt{\hat v_t}}$ 会逐坐标扭曲正则化项：
\[
  \Delta w_t = -\eta\frac{\grad\mathcal{L}(w_t) + \lambda w_t}{\sqrt{\hat v_t}+\eps}
  \neq (1-\eta\lambda)w_t - \eta\frac{\grad\mathcal{L}(w_t)}{\sqrt{\hat v_t}+\eps}.
\]
对梯度大的坐标，$\sqrt{\hat v_t}$ 大，正则化被削弱；对梯度小的坐标，反而加强。

\vspace{0.2cm}
\textbf{AdamW 解耦方案}：
\[
  \boxed{w_{t+1} = (1-\eta\lambda)w_t - \eta\frac{\hat m_t}{\sqrt{\hat v_t}+\eps}}
\]
Weight decay 独立于自适应更新，两个机制不互相干扰。
\end{frame}

\begin{frame}{Gradient Clipping}
梯度裁剪防止梯度爆炸（尤其RNN/Transformer训练中）：

\vspace{0.2cm}
\textbf{按值裁剪 (Value Clipping)}：
\[
  g_i = \begin{cases}
    c & \text{if } g_i > c\\
    -c & \text{if } g_i < -c\\
    g_i & \text{otherwise}
  \end{cases}
\]

\vspace{0.2cm}
\textbf{按范数裁剪 (Norm Clipping)}：
\[
  g \leftarrow \begin{cases}
    g & \text{if } \norm{g} \le v\\
    v \cdot \frac{g}{\norm{g}} & \text{otherwise}
  \end{cases}
\]

\vspace{0.2cm}
\textbf{为什么有效？} SGD 在光滑假设下需要 $\eta\le 2/L$，但深度学习中的 $L$ 可能非常大且不均匀。
梯度裁剪相当于人为限制了有效 $L$，使更大的学习率成为可能。
\end{frame}

% ======================================================================
\section{Adam 与自适应梯度}
% ======================================================================

\begin{frame}{Adam 的完整推导回顾}
Adam 维护两个 EMA：
\[
  \begin{aligned}
    m_t &= \beta_1 m_{t-1} + (1-\beta_1)g_t &\quad& \text{（一阶矩：方向平滑）}\\
    v_t &= \beta_2 v_{t-1} + (1-\beta_2)g_t^{\odot2} &\quad& \text{（二阶矩：逐坐标方差）}
  \end{aligned}
\]

\vspace{0.2cm}
\textbf{偏差修正}（关键但常被忽略的推导）：
\[
  \E[m_t] = (1-\beta_1)\sum_{k=1}^t \beta_1^{t-k}\E[g_k]
  = (1-\beta_1^t)\mu \;\Longrightarrow\; \hat m_t = \frac{m_t}{1-\beta_1^t}.
\]
同理 $\hat v_t = v_t/(1-\beta_2^t)$。修正后 $\E[\hat m_t]=\mu$，消除了冷启动偏差。

\vspace{0.2cm}
\textbf{更新}：
\[
  w_{t+1} = w_t - \eta\frac{\hat m_t}{\sqrt{\hat v_t}+\eps},\qquad
  \text{逐坐标操作：}\;\frac{\hat m_{t,i}}{\sqrt{\hat v_{t,i}}+\eps}.
\]

\vspace{0.2cm}
\textbf{预条件视角}：$w_{t+1}=w_t-\eta P_t\hat m_t$，其中 $P_t=\diag(\hat v_t+\eps^2)^{-1/2}$ 是对角矩阵。
\end{frame}

\begin{frame}{Adam 的 reinterpretation：Gauss-Newton 联系}
\textbf{关键洞察}：Adam 的 $\sqrt{v_t}$ 可以看作 Gauss-Newton 矩阵对角元的近似。

\vspace{0.2cm}
对于损失 $\ell(f(x;w),y)$，广义 Gauss-Newton (GGN) 矩阵为：
\[
  G(w) = J(w)^\top H_\ell(w) J(w),
\]
其中 $J(w)=\frac{\partial f}{\partial w}$ 是网络输出的 Jacobian，$H_\ell$ 是损失关于输出的 Hessian。

\vspace{0.2cm}
GGN 是 Hessian 的正定近似（忽略残差项），其对角元恰好是：
\[
  [G(w)]_{ii} = \sum_{j}\left(\frac{\partial f_j}{\partial w_i}\right)^2 \cdot [H_\ell]_{jj}.
\]
在特定条件下，$v_t$ 的 EMA 近似于 $\E[g_t^{\odot2}]$，而这正是 Fisher 信息矩阵（GGN 的期望）的对角元。

\vspace{0.2cm}
\textbf{所以 Adam 本质上在用对角 Gauss-Newton 做预条件} —— 它比「对角 Hessian」更准确地捕捉了深度网络的几何。
\end{frame}

\begin{frame}{Adam 收敛理论中的陷阱}
\textbf{原始证明的问题}（Reddi et al., 2018）：
Adam 的收敛需要有效学习率 $\eta_t = \eta/\sqrt{\hat v_t}$ 基本非增，但原始 Adam 中 $v_t$ 可能下降（$\beta_2$ 使旧信息衰减），导致 $\eta_t$ 在某些坐标上增大，破坏收敛。

\vspace{0.2cm}
\textbf{修正：AMSGrad}
\[
  \hat v_t = \max(\hat v_{t-1}, v_t),\qquad
  w_{t+1} = w_t - \frac{\eta}{\sqrt{t}}\frac{\hat m_t}{\sqrt{\hat v_t}+\eps}.
\]

\vspace{0.2cm}
\textbf{在线凸优化视角下 Adam 的 regret}：
在线预条件次梯度法的 regret bound（抽象形式）：
\[
  \text{Regret}_T \le \frac{1}{2\eta}\sum_{t=1}^T\big(\norm{\Delta_t}_{H_t}^2 - \norm{\Delta_{t+1}}_{H_t}^2\big)
  + \frac{\eta}{2}\sum_{t=1}^T g_t^\top H_t^{-1} g_t.
\]
取 $H_t=\diag(\sum_{s\le t}g_s^{\odot2})^{1/2}$ 即恢复 AdaGrad 的 regret bound；取对角 EMA 近似则对应 Adam。
\end{frame}

% ======================================================================
\section{Shampoo 与 Kronecker 预条件}
% ======================================================================

\begin{frame}{从向量到矩阵：为什么需要结构化的预条件}
Full-matrix AdaGrad 维护 $S_t = \sum_{s\le t} g_s g_s^\top \in \R^{mn\times mn}$，
存储代价 $\mathcal{O}(m^2 n^2)$ 不可行。

\vspace{0.2cm}
\textbf{观察}：对矩阵参数 $W\in\R^{m\times n}$，梯度 $G_t$ 的外积 $\vect(G_t)\vect(G_t)^\top$ 具有近似的 Kronecker 结构。

\vspace{0.2cm}
\textbf{关键不等式}（Shampoo 的理论基础）：
若 $\rank(G)=r$，则：
\[
  \frac{1}{r} g g^\top \preceq I_m \otimes (G^\top G),\qquad
  \frac{1}{r} g g^\top \preceq (G G^\top) \otimes I_n.
\]
更一般地：
\[
  \eps I_{mn} + \sum_{t} g_t g_t^\top \preceq
  \left(\eps I_m + \sum_t G_t G_t^\top\right)^{1/2}
  \otimes
  \left(\eps I_n + \sum_t G_t^\top G_t\right)^{1/2}
  = L_t^{1/2} \otimes R_t^{1/2}.
\]

\vspace{0.2cm}
\textbf{直观}：$L_t=G_t G_t^\top$ 捕获行空间相关性（输出神经元之间），$R_t=G_t^\top G_t$ 捕获列空间相关性（输入特征之间）。
\end{frame}

\begin{frame}{Shampoo 算法}
\textbf{状态维护}（每个矩阵参数）：
\[
  L_t = L_{t-1} + G_t G_t^\top \in \R^{m\times m},\qquad
  R_t = R_{t-1} + G_t^\top G_t \in \R^{n\times n}.
\]

\vspace{0.2cm}
\textbf{更新公式}：
\[
  \boxed{W_{t+1} = W_t - \eta\,L_t^{-1/4}\,G_t\,R_t^{-1/4}}
\]

\vspace{0.2cm}
\textbf{为什么是 \(-1/4\) 次幂？}
在向量化空间中，预条件器为：
\[
  \vect(W_{t+1}) = \vect(W_t) - \eta\,(R_t^{-1/4}\otimes L_t^{-1/4})\,\vect(G_t).
\]
若 $H_t \approx L_t^{1/2}\otimes R_t^{1/2}$，则 $H_t^{-1/2}\approx L_t^{-1/4}\otimes R_t^{-1/4}$。

\vspace{0.2cm}
内存对比：$\underbrace{mn}_{\text{Adam}}\;<\;\underbrace{m^2+n^2}_{\text{Shampoo}}\;\ll\;\underbrace{m^2 n^2}_{\text{full-matrix}}$
\end{frame}

\begin{frame}{Kronecker 积的性质速查（知识补充）}
\begin{definition}[Kronecker 积]
  $A\otimes B = [a_{ij}B]_{ij}\in\R^{mp\times nq}$，其中 $A\in\R^{m\times n},B\in\R^{p\times q}$。
\end{definition}

\vspace{0.15cm}
\textbf{需熟记的四条性质}：
\begin{enumerate}
  \item \textbf{混合积}：$(A\otimes B)(C\otimes D) = AC\otimes BD$（维度匹配时）
  \item \textbf{逆}：$(A\otimes B)^{-1}=A^{-1}\otimes B^{-1}$（$A,B$ 可逆）
  \item \textbf{幂函数}：对 $A,B\succ0$，$(A\otimes B)^\alpha = A^\alpha\otimes B^\alpha$
  \item \textbf{向量化}：$\vect(AXB) = (B^\top\otimes A)\vect(X)$
\end{enumerate}

\vspace{0.15cm}
\textbf{性质 3 的证明（谱分解法）}：
\[
  \begin{aligned}
    (A\otimes B)^\alpha &= (Q_A\Lambda_A Q_A^\top \otimes Q_B\Lambda_B Q_B^\top)^\alpha\\
    &= [(Q_A\otimes Q_B)(\Lambda_A\otimes\Lambda_B)(Q_A\otimes Q_B)^\top]^\alpha\\
    &= (Q_A\otimes Q_B)(\Lambda_A^\alpha\otimes\Lambda_B^\alpha)(Q_A\otimes Q_B)^\top
    = A^\alpha\otimes B^\alpha.
  \end{aligned}
\]

\vspace{0.15cm}
\textbf{性质 4 在 Shampoo 中的应用}：
\[
  \vect(L^{-1/4}GR^{-1/4}) = (R^{-1/4}\otimes L^{-1/4})\vect(G).
\]
\end{frame}

% ======================================================================
\section{SOAP：在特征基中运行 Adam}
% ======================================================================

\begin{frame}{SOAP：Shampoo 的特征基 + Adam 的 EMA}
\textbf{核心动机}：
\begin{itemize}
  \item Shampoo 给出好几何（$L_t,R_t$ 的特征方向），但缺乏 EMA 的平滑
  \item Adam 有 EMA 的稳定性，但只能做对角预条件（坐标轴受限）
  \item SOAP = 在 Shampoo 特征基中跑 Adam
\end{itemize}

\vspace{0.2cm}
\textbf{算法}：
\begin{enumerate}
  \item 维护 Shampoo 因子 $L_t,R_t$，每 $K$ 步做 eigendecomposition 得到 $Q_L,Q_R$
  \item 旋转梯度：$\tilde G_t = Q_L^\top G_t Q_R$
  \item 在 $\tilde G_t$ 坐标中运行 Adam：$\tilde m_t,\tilde v_t$（EMA）
  \item 更新：$W_{t+1}=W_t-\eta\,Q_L\!\left(\frac{\hat{\tilde m}_t}{\sqrt{\hat{\tilde v}_t}+\eps}\right)\!Q_R^\top$
\end{enumerate}

\vspace{0.2cm}
\textbf{关键工程细节}：当 $Q_L,Q_R$ 更新时，Adam 状态需变换到新基：
\[
  \tilde m_t^{\text{new}} = (Q_L^{\text{new}})^\top Q_L^{\text{old}}\;\tilde m_t^{\text{old}}\;(Q_R^{\text{old}})^\top Q_R^{\text{new}}.
\]
不这么做的话，$\tilde m_t$ 会对应错误的方向——相当于 Adam 突然"失忆"。
\end{frame}

% ======================================================================
\section{Muon：谱正交化动量}
% ======================================================================

\begin{frame}{Muon：从提出到轰动}
\textbf{Muon = Momentum Orthogonalized by Newton-Schulz}

\vspace{0.2cm}
由 Keller Jordan、Jeremy Bernstein 等人在 2024 年底提出，迅速成为该年度最受关注的优化器之一。

\vspace{0.2cm}
\textbf{核心算法（极简）}：
\[
  \begin{aligned}
    B_t &= \mu B_{t-1} + G_t \qquad\text{（动量累积）}\\
    O_t &= \text{NewtonSchulz}_5(B_t) \qquad\text{（正交化）}\\
    W_{t+1} &= W_t - \eta\,O_t
  \end{aligned}
\]

\vspace{0.2cm}
\textbf{惊人效果}：
\begin{itemize}
  \item NanoGPT speedrunning：训练速度提升 $1.35\times$
  \item 1.5B 参数 Transformer 达到 GPT-2 XL 水平：Muon 用 10 GPU-hours，AdamW 需 13.3 小时
  \item 仅需一阶动量（无二阶矩），\textbf{内存比 AdamW 更轻}
\end{itemize}
\end{frame}

\begin{frame}{Newton-Schulz 迭代：矩阵正交化的高效引擎}
\textbf{目标}：给定矩阵 $M$，求 $\argmin_{\|O\|_{op}\le 1} \|O-M\|_F$（最优解为矩阵符号函数）。

\vspace{0.15cm}
\textbf{SVD 视角}：若 $M=U\Sigma V^\top$，则最优 $O^\star=UV^\top$（将奇异值全部置 1）。

\vspace{0.15cm}
\textbf{Newton-Schulz 迭代（5 步版本）}：
\begin{enumerate}
  \item 归一化：$X\leftarrow X/(\|X\|_F+\eps)$
  \item 若行数 > 列数，则转置：$X\leftarrow X^\top$
  \item 迭代 5 次（多项式系数经精心设计）：
  \[
    \begin{aligned}
      A &= X X^\top\\
      B &= 3.4445A - 4.7750A^2 + 2.0315A^3\\
      X &\leftarrow 3.4445X + BX
    \end{aligned}
  \]
  \item 若曾转置，则转回
\end{enumerate}

\vspace{0.15cm}
\textbf{为什么用多项式迭代而非 SVD？}
\begin{itemize}
  \item SVD 在 bfloat16 下数值不稳定；多项式迭代在 bfloat16 下表现好得多
  \item 多项式迭代全是矩阵乘法，GPU 友好（适合 CUDA/Tensor Cores）
\end{itemize}
\end{frame}

\begin{frame}{Newton-Schulz 迭代的数学原理（详解）}
设 $G=U\Sigma V^\top$。一步 NS 迭代产生：
\[
  \begin{aligned}
    G' &= (a I + b GG^\top + c (GG^\top)^2) G\\
    &= (a UU^\top + b U\Sigma^2 U^\top + c U\Sigma^4 U^\top) U\Sigma V^\top\\
    &= U\,(a\Sigma + b\Sigma^3 + c\Sigma^5)\,V^\top.
  \end{aligned}
\]

\vspace{0.15cm}
定义多项式 $\phi(x) = ax + bx^3 + cx^5$。每步 NS 迭代对奇异值逐元素作用 $\phi$：
\[
  \phi_N(\sigma) = \underbrace{\phi\circ\phi\circ\cdots\circ\phi}_{N\text{ 次}}(\sigma).
\]

\vspace{0.15cm}
\textbf{系数设计目标}：使 $\phi_N(\sigma)\to 1$（对所有 $\sigma\in[0,1]$）。

\vspace{0.15cm}
NS 使用的特定系数 $(3.4445, -4.7750, 2.0315)$ 是精心设计的 5 次多项式，使得：
\begin{itemize}
  \item 在 $[0,1]$ 上 $\phi(\sigma)$ 尽可能接近 1
  \item 收敛速度快（二次收敛到符号函数）
  \item 在 bfloat16 精度下稳定
\end{itemize}
\end{frame}

\begin{frame}{Muon 的直观：为什么正交化有帮助？}
\textbf{Keller Jordan 的观察}：
\begin{quote}
\small
「对 Transformer 中 2D 参数的更新矩阵进行人工检查后发现，SGD-momentum 和 Adam 产生的更新\textbf{条件数极高}。
它们几乎是低秩矩阵——所有神经元的更新由极少数方向主导。正交化有效地提升了那些幅度小但对学习重要的『稀有方向』的尺度。」
\end{quote}

\vspace{0.2cm}
\textbf{数学解释}：
\begin{itemize}
  \item 梯度 $G_t$ 的奇异值分布极不均衡：少数奇异值很大，多数几乎为零
  \item 直接用 $G_t$ 更新 = 只在少数主导方向上有实质性进展
  \item 正交化 $G_t\mapsto UV^\top$：所有非零奇异值都被置为 1，每个方向获得平等权重
  \item 这相当于一种\textbf{谱预条件}：在奇异向量基中完全消除了奇异值的尺度差异
\end{itemize}
\end{frame}

% ======================================================================
\section{最陡下降统一视角}
% ======================================================================

\begin{frame}{最陡下降：不同范数 $\to$ 不同优化器}
\textbf{统一框架}：许多优化器可解释为不同范数下的最陡下降。
\[
  w_{t+1} = \argmin_w \left\{\langle g_t, w-w_t\rangle + \frac{1}{2\eta}\|w-w_t\|^2\right\}
\]
其中 $\|\cdot\|$ 是\textbf{任意范数}（不一定来自内积！）。

\vspace{0.2cm}
\textbf{一般解法}：分解 $\delta = c\cdot u$（$u$ 是单位向量），最优方向 $u^\star$ 是 $-g_t$ 在对偶范数单位球上的最远点。

\vspace{0.2cm}
\begin{center}
\begin{tabular}{lll}
\toprule
范数 $\|\cdot\|$ & 对偶范数 & 产生的优化器 \\
\midrule
$\ell_2$ & $\ell_2$ & 标准 SGD / GD \\
$\ell_\infty$ & $\ell_1$ & SignSGD（逐坐标取符号） \\
$\|\cdot\|_F$（Frobenius） & $\|\cdot\|_F$ & 标准 SGD（矩阵形式） \\
$\|\cdot\|_{op}$（谱范数） & $\|\cdot\|_1$（核范数） & Muon / PolarGrad \\
\bottomrule
\end{tabular}
\end{center}

\textbf{关键结论}：Muon = 在\textbf{谱范数}下的最陡下降方向（$UV^\top$）。
\end{frame}

\begin{frame}{谱范数与矩阵符号函数（知识补充）}
\begin{definition}[谱范数 / 算子范数]
  对矩阵 $G\in\R^{m\times n}$，谱范数（最大奇异值）：
  \[
    \|G\|_{op} = \sigma_{\max}(G) = \max_{\|x\|_2=1}\|Gx\|_2.
  \]
\end{definition}

\begin{definition}[核范数 / 迹范数]
  \[
    \|G\|_* = \sum_i \sigma_i(G) = \tr(\sqrt{G^\top G}).
  \]
\end{definition}

\vspace{0.15cm}
\textbf{核范数是谱范数的对偶范数}：$\|G\|_* = \sup_{\|X\|_{op}\le 1}\langle G,X\rangle$。

\vspace{0.15cm}
\begin{theorem}[谱范数下的最陡方向]
  \[
    \argmax_{\|U\|_{op}\le 1}\langle -G, U\rangle = -U_G V_G^\top = -\sgn(G),
  \]
  其中 $G=U_G\Sigma V_G^\top$ 是 SVD，$\sgn(G)$ 是矩阵符号函数。
\end{theorem}

\vspace{0.15cm}
\textbf{极分解视角}：$\sgn(G)$ = 极分解 $G=UP$ 中的正交因子 $U$（其中 $P\succeq0$）。
这是 Muon 的数学本质：\textbf{丢弃梯度的尺度信息，只保留方向（正交因子）}。
\end{frame}

% ======================================================================
\section{苏炜杰的 PolarGrad 框架}
% ======================================================================

\begin{frame}{苏炜杰 (Weijie Su) 的关键贡献：曲率 vs 梯度预条件}
苏炜杰 (UPenn) 在 2025 年提出了一个统一框架，精确区分了两种预条件策略：

\vspace{0.2cm}
\begin{columns}[T]
\column{0.48\textwidth}
\begin{block}{曲率预条件 (Curvature Preconditioning)}
  代表：\textbf{Adam, AdaGrad}
  \begin{itemize}
    \item 目标：应对 Hessian 的各向异性
    \item 方法：估计并缩放每个参数维度的曲率
    \item 代价：忽略参数的矩阵结构
    \item 问题：可能引入训练不稳定性
  \end{itemize}
\end{block}

\column{0.48\textwidth}
\begin{block}{梯度预条件 (Gradient Preconditioning)}
  代表：\textbf{Muon, PolarGrad}
  \begin{itemize}
    \item 目标：应对梯度本身的各向异性
    \item 方法：正交化/极分解梯度矩阵
    \item 优势：保持矩阵结构，更稳定
    \item 局限：不直接处理曲率差异
  \end{itemize}
\end{block}
\end{columns}

\vspace{0.2cm}
\begin{center}
\textbf{Muon 的成功来自梯度预条件；Adam 的优势来自曲率预条件。}\\
\textbf{PolarGrad 尝试结合两者。}
\end{center}
\end{frame}

\begin{frame}{PolarGrad：基于极分解的统一优化器}
\textbf{极分解 (Polar Decomposition)}：任何矩阵 $G$ 可唯一分解为
\[
  G = U P,\qquad U^\top U = I,\; P \succeq 0.
\]
其中 $U$ 是正交因子（方向），$P=(G^\top G)^{1/2}$ 是正半定因子（尺度）。

\vspace{0.2cm}
\textbf{PolarGrad 的核心更新}：
\[
  W_{t+1} = W_t - \eta\cdot U_t\cdot \phi(P_t),
\]
其中 $(U_t,P_t)$ 是动量矩阵的极分解，$\phi(\cdot)$ 是对 $P_t$ 的谱函数（对标量特征值逐元素作用）。

\vspace{0.2cm}
\textbf{特例}：
\begin{itemize}
  \item $\phi(P)=I$（忽略尺度）$\to$ \textbf{Muon}
  \item $\phi(P)=P$（保留原始尺度）$\to$ 动量 SGD
  \item $\phi(P)=\|G\|_*/\sqrt{d}\cdot I$（按核范数缩放）$\to$ \textbf{PolarGrad 推荐配置}
\end{itemize}

\vspace{0.2cm}
\textbf{PolarGrad 的优势}：使用 QDWH 或 ZOLO-PD 等稳定的极分解算法（相比 NS 需要更少调参），在 LLM 预训练中全面超越 Adam、Shampoo 和 Muon。
\end{frame}

\begin{frame}{ASGO：单侧预条件器（讲座内容延续）}
\textbf{ASGO (Adaptive Structured Gradient Optimization)} 只用一个预条件矩阵（而非 Shampoo 的两个）：

\vspace{0.15cm}
\textbf{更新}：
\[
  W_{t+1} = W_t - \eta_t \Lambda_t^{-1} G_t,\qquad
  \Lambda_t = V_t^{1/2} + \eps I_m,\qquad
  V_t = V_{t-1} + G_t G_t^\top.
\]

\vspace{0.15cm}
\textbf{收敛保证}（非光滑凸）：
\[
  \frac{1}{T}\sum_{t=0}^{T-1}\E[f(W_t)] - f(W^\star) \le \mathcal{O}\!\left(\frac{\|Q\|_* D_{op}}{\sqrt{T}}\right)
\]
其中 $\|Q\|_*$ 是梯度协方差上界的核范数，$D_{op}$ 是权重差异的谱范数。

\vspace{0.15cm}
\textbf{与 Muon 的关系}：
\[
  \text{ASGO} : \text{Muon} \;=\; \text{AdaGrad} : \text{SignSGD}
\]
ASGO 在谱范数最陡下降框架中加入自适应预条件（核范数缩放），Muon 只用固定的正交化。

\vspace{0.15cm}
\textbf{实验验证}：GPT-2 (124M) 预训练中，ASGO-PE 的验证 loss (3.342) 与 Muon (3.342) 持平，均优于 AdamW (3.455) 和 Shampoo (3.417，不稳定)。
\end{frame}

% ======================================================================
\section{完整数学工具箱}
% ======================================================================

\begin{frame}{你需要掌握的线性代数工具}
\begin{enumerate}
  \item \textbf{谱分解}：$A=Q\Lambda Q^\top$（对称矩阵）、$\Lambda=\diag(\lambda_1,\ldots,\lambda_d)$
  \item \textbf{奇异值分解 (SVD)}：$G=U\Sigma V^\top$（任意矩阵）、$\Sigma=\diag(\sigma_1,\ldots,\sigma_r)$
  \item \textbf{极分解}：$G=UP$，$U$ 正交，$P\succeq0$。关系：$P=V\Sigma V^\top$，$U=UV^\top$（SVD 的 $U,V$）
  \item \textbf{矩阵函数}：$\phi(A)=Q\phi(\Lambda)Q^\top$（对对称矩阵）。\textbf{注意}：$\phi(G)$ 不直接对非对称矩阵定义 \\
  $\to$ 需通过 SVD：$\phi_{\text{sing}}(G)=U\phi(\Sigma)V^\top$
  \item \textbf{矩阵符号函数}：$\sgn(G)=UV^\top$（将奇异值全部置 1）
  \item \textbf{三类矩阵范数}：
  \begin{itemize}
    \item Frobenius：$\|G\|_F = \sqrt{\sum_{i,j}G_{ij}^2} = \sqrt{\sum_i\sigma_i^2}$
    \item 谱范数：$\|G\|_{op} = \sigma_{\max}$
    \item 核范数：$\|G\|_* = \sum_i\sigma_i$
  \end{itemize}
  \item \textbf{Loewner 偏序}：$A\preceq B \iff B-A\succeq0$（$B-A$ 正半定）
\end{enumerate}
\end{frame}

\begin{frame}{极分解 vs SVD vs 谱分解的辨析（易错点）}
\begin{center}
\begin{tabular}{p{2.5cm}p{3.5cm}p{3.5cm}p{3cm}}
\toprule
 & \textbf{谱分解} & \textbf{SVD} & \textbf{极分解} \\
\midrule
适用矩阵 & 对称矩阵 $A=A^\top$ & 任意矩阵 $G$ & 任意方阵/矩阵 \\
形式 & $A=Q\Lambda Q^\top$ & $G=U\Sigma V^\top$ & $G=UP$ \\
因子 & $Q$ 正交，$\Lambda$ 对角 & $U,V$ 正交，$\Sigma$ 非负对角 & $U$ 正交，$P\succeq0$ \\
关系 & SVD 的特例 & $\to$ 极分解：$U_{\text{polar}}=UV^\top$ & $P=V\Sigma V^\top$ \\
\bottomrule
\end{tabular}
\end{center}

\vspace{0.3cm}
\textbf{常见错误}：
\begin{itemize}
  \item 对非对称矩阵做"谱分解" → 不存在！请用 SVD。
  \item 混淆 SVD 的 $U$ 与极分解的 $U$ → SVD 有两个正交矩阵，极分解的 $U=U_{\text{SVD}}V_{\text{SVD}}^\top$。
  \item $\phi(G)$ 对非对称矩阵没有标准定义 → 若论文中出现，通常指对奇异值作用 $\phi$。
\end{itemize}
\end{frame}

\begin{frame}{矩阵函数的定义与性质（知识补充）}
\begin{definition}[对称矩阵的函数]
  若 $A=Q\Lambda Q^\top$（$\Lambda=\diag(\lambda_1,\ldots,\lambda_d)$），则
  \[
    \phi(A) := Q\,\diag(\phi(\lambda_1),\ldots,\phi(\lambda_d))\,Q^\top.
  \]
\end{definition}

\vspace{0.15cm}
\textbf{为什么需要这个定义？} 矩阵函数不是逐元素函数！
\[
  \exp(A) \neq [e^{a_{ij}}],\qquad
  A^{1/2} \neq [\sqrt{a_{ij}}],\qquad
  \log(A) \neq [\log a_{ij}].
\]
逐元素操作（如 Adam 的 $g_t^{\odot2}$）是 Hadamard 代数，与矩阵函数是两回事。

\vspace{0.15cm}
\textbf{保证矩阵函数良定义的条件}：
\begin{itemize}
  \item $A$ 可对角化（对称矩阵自动满足）
  \item $\phi$ 在 $A$ 的谱上有定义（如 $\phi(x)=x^{-1/2}$ 要求所有 $\lambda_i>0$）
  \item $\to$ 实践中用 damping $A+\eps I$ 确保正定性
\end{itemize}
\end{frame}

\begin{frame}{Loewner 偏序与矩阵不等式（知识补充）}
\begin{definition}[Loewner 偏序]
  $A \preceq B \iff B-A$ 是正半定矩阵（特征值全 $\ge0$）。
\end{definition}

\vspace{0.2cm}
\textbf{关键性质}：
\begin{itemize}
  \item 若 $A\preceq B$，则对任意 $X$，$X^\top A X \preceq X^\top B X$
  \item 若 $A\preceq B$ 且 $A,B\succ0$，则 $B^{-1}\preceq A^{-1}$（序反转）
  \item 对任意函数 $\phi$，$\phi(A)\preceq\phi(B)$ \textbf{不保证}成立（需 $\phi$ 是算子单调的）
  \item $\phi(x)=x^\alpha$ 对 $\alpha\in[0,1]$ 是算子单调的 $\to$ $A^\alpha\preceq B^\alpha$ 当 $A\preceq B$
\end{itemize}

\vspace{0.2cm}
\textbf{Shampoo 中的不等式}：
\[
  g_t g_t^\top \preceq r\cdot (I_m\otimes (G_t^\top G_t))
\]
意思是：向量化梯度的外积在 Loewner 序下被 Kronecker 结构的上界控制。
这保证了 Shampoo 的预条件器不会"过度自信"（regret bound 的证明依赖此不等式）。
\end{frame}

% ======================================================================
\section{实践对比与建议}
% ======================================================================

\begin{frame}{各优化器对比总结}
\begin{center}
\scriptsize
\begin{tabular}{p{1.6cm}p{2.2cm}p{2.2cm}p{2.2cm}p{2.2cm}p{2cm}}
\toprule
维度 & SGD+Momentum & Adam/AdamW & Shampoo & Muon & PolarGrad \\
\midrule
预条件 & 无（或固定 LR） & 对角自适应 & Kronecker 结构 & 谱正交化 & 极分解 \\
额外内存 & $d$（动量） & $2d$（$m,v$） & $m^2+n^2+d$ & $d$（仅动量） & $\approx2d$ \\
矩阵感知 & 否 & 否 & 是 & 是 & 是 \\
曲率适应 & 否 & 是（对角） & 是（结构） & 否 & 部分 \\
需要二阶矩 & 否 & 是 & 否（但有外积累积） & 否 & 否 \\
矩阵根 & 否 & 否 & 是（$-1/4$ 次） & 否（NS 迭代） & 否（极分解） \\
2025 性能 & 基准 & 标准默认 & 中等规模好 & SOTA 速度 & SOTA 综合 \\
\bottomrule
\end{tabular}
\end{center}

\vspace{0.3cm}
\textbf{实践建议}（2025-2026 状态）：
\begin{itemize}
  \item \textbf{默认起点}：AdamW（最成熟、框架支持最好）
  \item \textbf{追求训练速度}：Muon（内存更轻、收敛更快）
  \item \textbf{追求最佳性能}：PolarGrad（理论上更完备，但生态不如 AdamW 成熟）
  \item \textbf{大 batch 训练}：Shampoo / SOAP
\end{itemize}
\end{frame}

% ======================================================================
\section{练习}
% ======================================================================

\begin{frame}{理论练习}
\begin{enumerate}
  \item \textbf{SAA 收敛}：对有限域 $|\mathcal{X}|=K$，推导 SAA 的样本复杂度 bound。说明为什么结果是 $\mathcal{O}(\log K)$ 而非 $\mathcal{O}(K)$（这体现了联合界不是紧的，实际上可以用覆盖数得到更好的 bound）。

  \item \textbf{Hoeffding 的证明}：从 Chernoff bound $\Pr(Z_N\ge a)\le e^{-ta}\E[e^{t Z_N}]$ 出发，对 sub-Gaussian $Y_i$ 证明 $\Pr(\frac1N\sum_i Y_i\ge a)\le e^{-Na^2/(2\sigma^2)}$。关键在于优化 $t$ 的选择。

  \item \textbf{Mirror Descent 的三点恒等式}：证明对 Bregman divergence $D_\psi$，
  \[
    D_\psi(z,x) - D_\psi(z,y) - D_\psi(y,x) = \langle\grad\psi(y)-\grad\psi(x),z-y\rangle.
  \]
  并解释为什么这个恒等式是 regret 分析的核心工具。

  \item \textbf{谱范数最陡方向}：设 $G=U\Sigma V^\top$，求解 $\argmin_{\|O\|_{op}\le 1}\|O-G\|_F^2$。证明最优解为 $O^\star=UV^\top$。（提示：利用 von Neumann 迹不等式。）
\end{enumerate}
\end{frame}

\begin{frame}{工程练习}
\begin{enumerate}
  \item \textbf{实现 Muon 的简化版}：在二维矩阵优化问题上（如矩阵分解），实现带动量的 Newton-Schulz（2 次迭代即可），可视化更新矩阵的奇异值演化。

  \item \textbf{Adam vs Muon 的旋转敏感性}：构造旋转二次矩阵函数 $f(W)=\frac12\|AW-W^\star\|_F^2$，分别用 Adam 和 Muon 优化。改变 $A$ 的条件数，记录两者的收敛步数。

  \item \textbf{Newton-Schulz 收敛实验}：随机生成 $100\times100$ 矩阵 $M$，计算 $\sgn(M)=UV^\top$（通过 SVD 作为 ground truth）。实现 NS 迭代并记录 $\|X_k-\sgn(M)\|_F$ 随 $k$ 的变化（1-10 次迭代）。验证二次收敛性。

  \item \textbf{可视化 Muon 的预条件效果}：对两层 MLP 在 MNIST 上训练，每 10 步记录梯度矩阵 $G_t$ 的奇异值分布（直方图）。对比 AdamW 和 Muon 下奇异值分布的变化。
\end{enumerate}
\end{frame}

% ======================================================================
\section{参考文献}
% ======================================================================

\begin{frame}{核心参考文献}
\footnotesize
\begin{thebibliography}{99}

\bibitem{shapiro}
A. Shapiro, D. Dentcheva, A. Ruszczy\'nski.
\newblock \textit{Lectures on Stochastic Programming: Modeling and Theory}, SIAM, 2021.

\bibitem{adam}
D. P. Kingma \& J. Ba.
\newblock Adam: A Method for Stochastic Optimization. \textit{ICLR}, 2015.

\bibitem{adamw}
I. Loshchilov \& F. Hutter.
\newblock Decoupled Weight Decay Regularization. \textit{ICLR}, 2019.

\bibitem{amsgrad}
S. J. Reddi, S. Kale, \& S. Kumar.
\newblock On the Convergence of Adam and Beyond. \textit{ICLR}, 2018.

\bibitem{shampoo}
V. Gupta, T. Koren, \& Y. Singer.
\newblock Shampoo: Preconditioned Stochastic Tensor Optimization. \textit{ICML}, 2018.

\bibitem{soap}
N. Vyas et al.
\newblock SOAP: Improving and Stabilizing Shampoo using Adam. \textit{arXiv:2409.11321}, 2024.

\bibitem{muon}
K. Jordan, J. Bernstein et al.
\newblock Muon: An Optimizer for Matrix Structured Gradient. \textit{arXiv}, 2024.

\bibitem{polar}
T. T.-K. Lau, Q. Long, \& W. J. Su.
\newblock PolarGrad: A Class of Matrix-Gradient Optimizers from a Unifying Preconditioning Perspective. \textit{arXiv:2505.21799}, 2025.

\bibitem{isotropic}
W. J. Su et al.
\newblock Isotropic Curvature Model for Understanding Deep Learning Optimization. \textit{arXiv:2511.00674}, 2025.

\bibitem{asgo}
X. Chen et al.
\newblock ASGO: Adaptive Structured Gradient Optimization. 2025.

\bibitem{steepest}
J. Bernstein \& L. Newhouse.
\newblock A Unifying Framework for Optimization Algorithms. \textit{arXiv}, 2024.

\end{thebibliography}
\end{frame}

\begin{frame}{结束页}
\centering
\Large
\textbf{核心要点回顾}

\vspace{0.3cm}
\normalsize
\begin{enumerate}
  \item 算法复杂度 = SAA（样本多少够？）+ SA/SGD（迭代多少步够？）
  \item AdamW 的关键创新：weight decay 与自适应梯度解耦
  \item Muon 的数学本质：谱范数下的最陡下降 = 梯度矩阵的极分解（正交因子）
  \item 苏炜杰 (Weijie Su) 的框架：区分「曲率预条件」与「梯度预条件」
  \item PolarGrad 统一了 Muon/Shampoo/Adam 在极分解视角下
\end{enumerate}

\vspace{0.4cm}
下一步：用二维 toy problem 亲手验证「Adam 怕旋转而 Muon 不怕」。
\end{frame}

\end{document}
